\section{Institutional setting: Indonesian FX fixing, taxation, and exchange costs}
\label{sec:institutional}

Understanding Indonesian cryptocurrency pricing requires tracking three institutional structures that govern local execution costs and benchmark timing: the Bank Indonesia JISDOR fixing calendar, statutory crypto tax regimes, and exchange plus clearing-house fee schedules.

\subsection{Bank Indonesia JISDOR publication regimes}
JISDOR is Bank Indonesia's official interbank USD/IDR reference rate calculated from interbank spot transactions reported through the Sistem Monitoring Transaksi Valuta Asing Terhadap Rupiah (SISMONTAVAR). Prior to April 2021, JISDOR observed morning transactions and published at 10:00 WIB (03:00 UTC). Under Bank Indonesia Press Release No.\ 23/30/DKom (issued February 3, 2021, effective April 5, 2021), Bank Indonesia expanded the underlying transaction window to 08:00--16:00 WIB and shifted publication to 16:15 WIB (09:15 UTC) \citep{bi_sp233021}. During pandemic operating adjustments, temporary 15:15 WIB (08:15 UTC) releases also occurred before Board of Governors Regulation PADG No.\ 7 Tahun 2026 formally moved daily publication to 15:15 WIB (08:15 UTC) effective April 1, 2026 \citep{bi_padg072026}.

Two methodological consequences follow from JISDOR's structure. First, JISDOR publishes strictly on Indonesian banking business days, leaving weekends and national holidays without an official fixing even though cryptocurrency spot books operate continuously. Forward-filling Friday JISDOR across weekends induces artificial predictability when FX markets move prior to Monday opening. Second, aligning daily 23:59 UTC cryptocurrency closes with a 08:15 UTC or 09:15 UTC JISDOR benchmark introduces a 14.75-to-15.75 hour timing mismatch. As Section~\ref{sec:robustness} verifies, this mismatch inflates pure-crypto basis dispersion from $27.44\text{ bps}$ to $35.97\text{ bps}$.

\subsection{Statutory tax and clearing regimes (2021--2026)}
Indonesian digital-asset transactions experienced six distinct cost regimes between 2021 and late 2026 (summarized in Table~\ref{tab:fees}, Panel~A):
\begin{enumerate}
  \item \textbf{Pre-Tax Baseline (January 2021--April 2022, R1):} No special crypto transaction VAT or income tax applied. Indodax charged $0\%$ maker and $0.30\%$ taker fees on IDR spot pairs, generating a $60.0\text{ bps}$ round-trip taker cost on a domestic two-leg trade (excluding on-chain settlement).
  \item \textbf{PMK No.\ 68/PMK.03/2022 (May 2022--October 2024, R2A--R2B):} Effective May 1, 2022, the Ministry of Finance imposed a final Value Added Tax (PPN) of $0.11\%$ on buyers ($0.22\%$ on unregistered platforms) and a final Article 22 Income Tax (PPh 22) of $0.10\%$ on sellers ($0.20\%$ unregistered) \citep{pmk68_2022}. Around August 22--23, 2023, Indodax revised its Pro trading fee schedule, cutting taker fees by roughly $40\%$ toward the modern $0.10\%$ maker / $0.20\%$ taker structure (VAT-inclusive on service fees). Because daily archive snapshots of the intermediate 2022--2023 USDT pair fee tier are only partially preserved publicly, regimes R2A and R2B carry an explicit \texttt{PARTIALLY\_VERIFIED} provenance flag and are bracketed with sensitivity bounds.
  \item \textbf{CFX and KKI Bourse Fee Launch (November 1, 2024--July 2025, R3--R3B):} Under BAPPEBTI Regulation PERBA No.\ 9/2024 and CFX/KKI Circulars SEB 003/004, trades executed on registered Indonesian venues incurred an additional bourse and clearing levy of $2.22\text{ bps}$ on IDR pairs and $4.44\text{ bps}$ on USDT pairs (inclusive of $11\%$ service VAT), effective 23:59:59 WIB on October 31, 2024. Although nominal VAT legislation adjusted in January 2025 (PMK 81/2024 and PMK 131/2024), the $11/12$ tax-base rule preserved the $0.11\%$ effective asset VAT rate. On January 10, 2025, primary market supervision formally transferred from BAPPEBTI to OJK under POJK No.\ 27/2024.
  \item \textbf{PMK No.\ 50 Tahun 2025 Overhaul (August 1, 2025--February 2026, R4):} Effective August 1, 2025, PMK No.\ 50 Tahun 2025 abolished buyer asset VAT ($0\text{ bps}$) while raising seller final PPh 22 from $10\text{ bps}$ to $21\text{ bps}$ on domestic licensed platforms ($100\text{ bps}$ on foreign platforms) \citep{pmk50_2025}. Statutory tax shifted completely onto the sell side, though symmetric two-leg round-trip friction remained $65.44\text{ bps}$ (taker) until bourse fee cuts.
  \item \textbf{Two-Stage CFX/KKI Halvings (March 2026 and October 2026, R5--R6):} Under SEB No.\ 001/2026 effective March 1, 2026, CFX/KKI halved levies to $1.11\text{ bps}$ (IDR) and $2.22\text{ bps}$ (USDT). A second $50\%$ reduction on October 1, 2026 brought CFX/KKI charges to $0.555\text{ bps}$ (IDR) and $1.11\text{ bps}$ (USDT). Summing the $10/20\text{ bps}$ Pro platform fee (VAT-inclusive), $0/21\text{ bps}$ statutory tax, and $0.555\text{ bps}$ CFX fee reproduces Indodax's live published Pro rates to the fifth decimal place ($0.10555\%/0.20555\%$ buy; $0.31555\%/0.41555\%$ sell; $0.2511\%/0.2811\%$ USDT).
\end{enumerate}

\begin{table}[htbp]
\centering
\caption{Indonesian statutory crypto tax, CFX/KKI, platform commission, and route execution costs (bps)}
\label{tab:fees}
\fontsize{6.3}{7.5}\selectfont\setlength{\tabcolsep}{1.8pt}
\begingroup
\begin{tabular}{@{}l>{\raggedright\arraybackslash}p{2.3cm}rrrrrr@{}}
\toprule
\multicolumn{8}{@{}l}{\textbf{Panel A: Chronological fee and tax regimes (2021--2026)}} \\
Window & Regime & M/T & VAT & PPh22 & CFX & Rte A (M/T) & Rte B (M/T) \\
\midrule
2021-01 to 2022-04 & R1: Pre-tax baseline & 0/30 & 0 & 0 & 0 & 0.0/60.0 & 7.5/70.0 \\
2022-05 to 2023-08 & R2A: PMK 68 (0/30)$^\dagger$ & 0/30 & 11 & 10 & 0 & 21.0/81.0 & 28.5/91.0 \\
2023-08 to 2024-10 & R2B: Indodax Pro cut$^\dagger$ & 10/20 & 11 & 10 & 0 & 41.0/61.0 & 48.5/71.0 \\
2024-11 to 2024-12 & R3: CFX fee launch & 10/20 & 11 & 10 & 2.22 & 45.4/65.4 & 52.9/75.4 \\
2025-01 to 2025-07 & R3B: PMK 81/131 VAT & 10/20 & 11 & 10 & 2.22 & 45.4/65.4 & 52.9/75.4 \\
2025-08 to 2026-02 & R4: PMK 50/2025 & 10/20 & 0 & 21 & 2.22 & 45.4/65.4 & 52.9/75.4 \\
2026-03 to 2026-09 & R5: CFX cut (Mar 2026) & 10/20 & 0 & 21 & 1.11 & 43.2/63.2 & 50.7/73.2 \\
2026-10 to 2026-11 & R6: CFX cut (Oct 2026) & 10/20 & 0 & 21 & 0.555 & 42.1/62.1 & 49.6/72.1 \\
\midrule
\multicolumn{8}{@{}l}{\textbf{Panel B: Live order-book costs by route and trade size (Nov 13, 2026 snapshot, R6)}} \\
Route & Scope & Size & \multicolumn{2}{r}{Stat. (Mkr/Tkr)} & \multicolumn{2}{r}{Depth Slip} & All-In Tkr \\
\midrule
Route A & Domestic 2-Leg & \$1k & \multicolumn{2}{r}{42.11/62.11} & \multicolumn{2}{r}{4.35} & 66.46 \\
Route B & Cross-Exchange 3-Leg & \$1k & \multicolumn{2}{r}{49.61/72.11} & \multicolumn{2}{r}{4.35} & 76.46 \\
Route C & Same-Exchange Indodax & \$1k & \multicolumn{2}{r}{67.22/90.22} & \multicolumn{2}{r}{63.14} & 153.36 \\
Route A & Domestic 2-Leg & \$10k & \multicolumn{2}{r}{42.11/62.11} & \multicolumn{2}{r}{4.35} & 66.46 \\
Route B & Cross-Exchange 3-Leg & \$10k & \multicolumn{2}{r}{49.61/72.11} & \multicolumn{2}{r}{4.35} & 76.46 \\
Route C & Same-Exchange Indodax & \$10k & \multicolumn{2}{r}{67.22/90.22} & \multicolumn{2}{r}{63.14} & 153.36 \\
Route A & Domestic 2-Leg & \$50k & \multicolumn{2}{r}{42.11/62.11} & \multicolumn{2}{r}{9.73} & 71.84 \\
Route B & Cross-Exchange 3-Leg & \$50k & \multicolumn{2}{r}{49.61/72.11} & \multicolumn{2}{r}{9.74} & 81.85 \\
Route C & Same-Exchange Indodax & \$50k & \multicolumn{2}{r}{67.22/90.22} & \multicolumn{2}{r}{68.53} & 158.75 \\
\bottomrule
\end{tabular}\endgroup
\begin{minipage}{\textwidth}\footnotesize\vspace{3pt}
\emph{Notes:} $^\dagger$The August 2023 platform fee transition (Regimes R2A--R2B) is only partially documented in public archives and is bracketed with sensitivity bounds in the empirical containment checks. Regime R6 matches published Indodax Pro rates to the fifth decimal place (0.10555\%/0.20555\% buy; 0.31555\%/0.41555\% sell; 0.2511\%/0.2811\% USDT). Route B and Route C calculations assume pre-funded inventory and exclude on-chain withdrawal fees.
\end{minipage}
\end{table}


\subsection{Route-specific execution architecture}
Using one cost threshold to evaluate both $B_t^{crypto}$ and $B_t^{total}$ confuses distinct execution pathways. Reconciling a domestic \texttt{BTC/IDR} quote against an offshore \texttt{BTC/USDT} quote requires three economic routes:
\textbf{Route A (Domestic Two-Leg Conversion)} sells \texttt{BTC/IDR} and buys \texttt{USDT/IDR} inside Indodax ($42.11\text{ bps}$ maker / $62.11\text{ bps}$ taker statutory fees in Regime R6, plus $4.35\text{ bps}$ live half-spread slippage at $\$10\text{k}$ notional for $66.46\text{ bps}$ all-in). \textbf{Route B (Cross-Exchange Three-Leg Arbitrage)} adds a Binance Spot \texttt{BTC/USDT} leg ($7.5\text{ bps}$ BNB / $10.0\text{ bps}$ standard fee, giving $72.11\text{ bps}$ statutory taker and $76.46\text{ bps}$ with $\$10\text{k}$ slippage, assuming pre-funded cross-venue inventory). \textbf{Route C (Same-Exchange Indodax \texttt{BTC/USDT})} avoids cross-exchange inventory transfer but incurs Indodax's USDT-pair fees and $4.44\text{ bps}$--$1.11\text{ bps}$ CFX USDT levies plus severe illiquidity ($117.59\text{ bps}$ live top-of-book spread, pushing Route C execution to $153.36\text{ bps}$--$158.75\text{ bps}$, Table~\ref{tab:fees} Panel~B). Conversely, none of these crypto routes settles straight into Bank Indonesia interbank USD/IDR without off-ramp banking fees, foreign-exchange spreads, or capital-conversion delays, explaining why $B_t^{stablecoin\text{-}FX}$ is structurally wider and more persistent than $B_t^{crypto}$.
